%%% BEGIN fig:nrqcd
\begin{figure}[htbp]
\centering
\begin{tikzpicture}[fd,
  t3 heavy/.style={draw,solid,line width=1.3pt,fd arrow=#1},
  t3 heavy/.default=0.5,
  t3 soft/.style={gluon,draw=green!45!black},
  t3 lab/.style={font=\small},
  t3 note/.style={font=\footnotesize}]
  % ===================== (a) colour-singlet channel =====================
  \begin{scope}
    \node[hard blob,circle,minimum size=1.0cm,inner sep=0pt] (H) at (-3.4,0) {};
    \node[soft blob,ellipse,minimum width=2.6cm,minimum height=1.3cm,inner sep=0pt] (L) at (0.8,0) {};
    \draw[gluon] (-5.0,0.85) -- (H.150);  \draw[gluon] (-5.0,-0.85) -- (H.210);
    \draw[gluon] (H.45) -- (-2.15,1.3);
    \draw[t3 heavy] (-2.92,0.17) -- (-0.4,0.17);
    \draw[t3 heavy] (-0.4,-0.17) -- (-2.92,-0.17);
    \draw[bound state] (L.east) -- (3.55,0);
    \node[soft blob,ellipse,minimum width=2.6cm,minimum height=1.3cm,inner sep=0pt] at (0.8,0) {};
    \node[hard blob,circle,minimum size=1.0cm,inner sep=0pt] at (-3.4,0) {};
    \node[t3 lab,left] at (-5.0,0.85) {$g$};  \node[t3 lab,left] at (-5.0,-0.85) {$g$};
    \node[t3 lab,right] at (-2.15,1.3) {$g$};
    \node[t3 lab] at (-1.38,0.5) {$Q\bar Q[{}^3S_1^{[1]}]$};
    \node[t3 note,align=center] at (0.8,0) {$\langle\mathcal O^{H}({}^3S_1^{[1]})\rangle$};
    \node[t3 lab,right] at (3.55,0) {$H\ (J/\psi)$};
    \node[t3 note] at (-3.4,-0.78) {$\sim1/m_Q$};
    \node[t3 note] at (0.8,-0.92) {$\sim1/(m_Qv)$};
    \node[t3 lab,anchor=west] at (-5.6,-1.45) {(a)\ 色单态：短距离过程须多放出一个硬胶子};
  \end{scope}
  % ===================== (b) colour-octet channel =====================
  \begin{scope}[yshift=-3.6cm]
    \node[hard blob,circle,minimum size=1.0cm,inner sep=0pt] (H) at (-3.4,0) {};
    \node[soft blob,ellipse,minimum width=2.6cm,minimum height=1.3cm,inner sep=0pt] (L) at (0.8,0) {};
    \draw[plain] (-5.0,0.85) -- (H.150);  \draw[plain] (-5.0,-0.85) -- (H.210);
    \draw[t3 heavy] (-2.92,0.17) -- (-0.4,0.17);
    \draw[t3 heavy] (-0.4,-0.17) -- (-2.92,-0.17);
    \draw[t3 soft] (1.25,0.55) -- (2.1,1.45);
    \draw[t3 soft] (1.25,-0.55) -- (2.1,-1.45);
    \draw[bound state] (L.east) -- (3.55,0);
    \node[soft blob,ellipse,minimum width=2.6cm,minimum height=1.3cm,inner sep=0pt] at (0.8,0) {};
    \node[hard blob,circle,minimum size=1.0cm,inner sep=0pt] at (-3.4,0) {};
    \node[t3 lab,left] at (-5.0,0.85) {$i$};  \node[t3 lab,left] at (-5.0,-0.85) {$j$};
    \node[t3 lab] at (-1.38,0.5) {$Q\bar Q[n^{[8]}]$};
    \node[t3 note,align=center] at (0.8,0) {$\langle\mathcal O^{H}(n^{[8]})\rangle$};
    \node[t3 lab,right] at (3.55,0) {$H$};
    \node[t3 note,right,text=green!35!black] at (2.1,1.45) {软胶子，$E\sim m_Qv^2$};
    \node[t3 note] at (-3.4,-0.78) {$\sim1/m_Q$};
    \node[t3 lab,anchor=west] at (-5.6,-1.75) {(b)\ 色八重态：颜色在长距离上由软胶子带走};
  \end{scope}
\end{tikzpicture}
\caption{重夸克偶素产生的 NRQCD 因子化图像。灰色团为短距离过程：在距离 $\sim1/m_Q$ 上产生处于确定色、自旋与轨道角动量态 $n={}^{2S+1}L_J^{[1,8]}$ 的 $Q\bar Q$ 对（粗线；箭头为费米子数），
其截面 $\dd\hat\sigma(ij\to Q\bar Q[n]+X)$ 可按 $\alpha_s(m_Q)$ 微扰计算；绿色团为长距离矩阵元 $\langle\mathcal O^{H}(n)\rangle$，描述 $Q\bar Q[n]$（相对速度 $v$，尺度 $\sim1/(m_Qv)$）演化为强子 $H$ 的概率，按 $v$ 的幂次计数（$c\bar c$：$v^2\approx0.3$；$b\bar b$：$v^2\approx0.1$）。
(a)~色单态：例如 $gg\to Q\bar Q[{}^3S_1^{[1]}]+g$，在 $v$ 的领头阶 $Q\bar Q$ 直接成为 $J/\psi$；为满足色与 $C$ 宇称守恒，短距离过程必须多放出一个硬胶子。
(b)~色八重态（如 ${}^1S_0^{[8]}$、${}^3S_1^{[8]}$、${}^3P_J^{[8]}$）：$Q\bar Q$ 对须放出能量 $\sim m_Qv^2$ 的软胶子才能成为色单态强子，相应矩阵元相对于 (a) 有额外的 $v$ 压低。
产生截面的因子化形式 $\dd\sigma(H)=\sum_n\dd\hat\sigma(Q\bar Q[n])\,\langle\mathcal O^H(n)\rangle$ 是猜想（[C]），对衰变的类似公式有全阶论证（[A]）；八重态算符须以 Wilson 线补全规范不变性（见正文）。}
\label{fig:nrqcd}
\end{figure}
%%% END fig:nrqcd
